\documentclass{article}
\usepackage{amsmath, amssymb, amsthm}
\usepackage{hyperref}
\usepackage{geometry}
\geometry{margin=1in}

\title{Erd\H{o}s Problem \#98}
\date{Last updated: 2025-08-31}
\author{Source: erdosproblems.com}

\begin{document}
\maketitle

\noindent\textbf{Status:} OPEN \quad \textbf{No prize}

\noindent\textbf{Tags:} geometry, distances

\medskip

\noindent\textbf{Problem Statement:}

Let $h(n)$ be such that any $n$ points in $\mathbb{R}^2$, with no three on a line and no four on a circle, determine at least $h(n)$ distinct distances. Does $h(n)/n\to \infty$?




Erd\H{o}s could not even prove $h(n)\geq n$. Pach has shown $h(n)&lt;n^{\log_23}$. Erd\H{o}s, F\"{u}redi, and Pach \cite{EFPR93} have improved this to\[h(n) &lt; n\exp(c\sqrt{\log n})\]for some constant $c&gt;0$.




References

[EFPR93] Erd\H{o}s, Paul and F\"{u}redi, Zolt\'{a}n and Pach, J\'{a}nos and Ruzsa, Imre Z., The grid revisited . Discrete Math. (1993), 189--196.

\noindent\textbf{Related OEIS sequences:} possible


\bigskip
\noindent\small{Source: \url{https://www.erdosproblems.com/98}}
\end{document}
